% Part:sets-functions-relations
% Chapter: size-of-sets
% Section: introduction

\documentclass[../../../include/open-logic-section]{subfiles}

\begin{document}

\olfileid[ar]{sfr}{siz}{int}

\olsection{مقدمة}

لما وضع جورج كانتور نظرية المجموعات في سبعينيات القرن التاسع عشر،
قصد في جملة ما قصد إلى تسويغ تصور اجتماع غير متناه؛ أي «لانهاية
بالفعل» بعبارة أهل العصور الوسطى. وكان بحث \emph{حجم} المجموعات
ركنًا من ذلك. فإذا تمايزت $\OLId{latin-ordinary}{a}$ و~$\OLId{latin-ordinary}{b}$ و~$\OLId{latin-ordinary}{c}$، ظهر للحدس أن
$\{\OLId{latin-ordinary}{a}, \OLId{latin-ordinary}{b}, \OLId{latin-ordinary}{c}\}$ \emph{أكبر} من~$\{\OLId{latin-ordinary}{a}, \OLId{latin-ordinary}{b}\}$. فهل المجموعات
اللانهائية كلها سواء في الحجم؟ سيتبين أن الأمر ليس كذلك.

وأول ما نحتاج إليه التعداد. فكل مجموعة منتهية تسرد بذكر
!!{element}sها كلها. وبعض غير المنتهيات كذلك يمكن سرد
!!{element}sها، على أن تكون القائمة نفسها غير منتهية. وتسمى هذه المجموعات !!{enumerable}. وقد أظهر كانتور النتيجة العجيبة: ليست كل
مجموعة لا نهائية !!{enumerable}. وبنهاية الفصل يتم لنا فهم برهانه.

\end{document}
